% =============================================================
% conclusion_template.tex
% One-page conclusion: restate, synthesize, limit, future work.
% Do NOT introduce new results.
% =============================================================

\section{Conclusion}\label{sec:conclusion}

This paper asks [QUESTION]. Exploiting [VARIATION], we find that
\headlineEstimate{} (\headlineSE{}) -- a [X]\% [increase / decrease]
relative to the sample mean. The result is robust to [LADDER] and
holds across [HETEROGENEITY DIMENSION].

% --- Synthesize, don't summarize -----------------------------
Two patterns stand out. First, [PATTERN 1] is consistent with
[INTERPRETATION], suggesting that [BROADER POINT]. Second,
[PATTERN 2] is harder to reconcile with [ALTERNATIVE INTERPRETATION],
implying that [BROADER POINT].

% --- Implications --------------------------------------------
For [POLICY DOMAIN], our findings imply [IMPLICATION 1]. The
magnitude is large enough to [DECISION-RELEVANT FRAMING].
For [THEORETICAL DEBATE], the evidence supports [SIDE OF DEBATE]
and challenges [OPPOSING VIEW].

% --- Honest limitations --------------------------------------
Several limitations warrant mention. First, our setting
[LIMITATION 1: e.g.\ involves a single state] limits external
validity to [SCOPE]. Second, we do not observe [LIMITATION 2:
e.g.\ post-2020 data]. Third, [LIMITATION 3: e.g.\ our identification
relies on assumptions that are most plausible in the short run].

% --- Future directions ---------------------------------------
Future research could [DIRECTION 1] and [DIRECTION 2]. Our
findings highlight the importance of [BROADER QUESTION], and we
hope they motivate further work on [TOPIC].
